\chapter{Tiền, nghề và tương lai còn mù mờ}
\markboth{Tiền, nghề và tương lai còn mù mờ}{Tiền, nghề và tương lai còn mù mờ}

\section{Nhà em cần tiền sớm, ước mơ để sau được không?}
\begin{truyen}{Nhà em cần tiền sớm, ước mơ để sau được không?}{Classroom Talk}
\chuhoa{M}{ột bạn hỏi rất đời: ``Nhà em không dư giả. Em nghĩ kiếm tiền sớm mới là chuyện chính. Ước mơ để sau được không thầy?''}

Thầy đáp: ``Được. Không phải ai cũng có đặc quyền mơ đẹp trước. Nhưng kể cả khi cần tiền sớm, em vẫn phải chọn cách kiếm tiền khiến mình lớn lên chứ không mòn đi.''

Bạn ấy hỏi lại: ``Nghĩa là sao?''

Thầy nói: ``Làm việc để sống là thật. Nhưng nếu chỉ nhìn cái trước mắt, em dễ bán rẻ những năm sau của mình.''

\textit{(When money is urgent, dreams may need to wait. But choices still matter because short-term survival can shape long-term limits.)}
\end{truyen}

\begin{baihoc}
Thực tế không giết ước mơ. Thực tế chỉ buộc mình chọn đường đi bền và tỉnh hơn.
\textit{(Reality does not kill dreams. It forces wiser and steadier paths toward them.)}
\end{baihoc}

\begin{ghinhoanh}
\textbf{Short English to remember:}

Urgency is real.

Do not let urgency steal your future.
\end{ghinhoanh}

\begin{tuhocnhanh}
1. Em đang cần tiền gấp đến mức nào? \textit{(How urgent is your financial reality?)}

2. Con đường nào vừa giúp em sống vừa giúp em tiến? \textit{(What path helps you survive and still move forward?)}
\end{tuhocnhanh}
\ngancach

\section{AI làm hết rồi, học sâu làm gì nữa}
\begin{truyen}{AI làm hết rồi, học sâu làm gì nữa}{Classroom Talk}
\chuhoa{C}{ó bạn nói nửa đùa nửa thật: ``Mai mốt AI làm hết, tụi em học cho lắm rồi cũng bị thay thôi. Vậy học sâu làm gì?''}

Thầy hỏi: ``Nếu máy làm phần dễ hơn cho em, vậy em muốn thành người chỉ bấm nút hay thành người biết đặt câu hỏi, kiểm tra và quyết định?''

Bạn ấy cười: ``Nghe vậy thì học vẫn chưa được nghỉ ha thầy.''

Thầy đáp: ``Đúng. Càng nhiều công cụ mạnh, người dùng càng phải bớt nông.''

\textit{(New tools do not remove the need to think. They raise the cost of being shallow and careless.)}
\end{truyen}

\begin{baihoc}
Công cụ thông minh không cứu nổi người lười nghĩ.
\textit{(Smart tools cannot rescue lazy thinking.)}
\end{baihoc}

\begin{ghinhoanh}
\textbf{Short English to remember:}

Tools grow stronger.

Thinking must grow deeper.
\end{ghinhoanh}

\begin{tuhocnhanh}
1. Em muốn trở thành người dùng công cụ hay phụ thuộc công cụ? \textit{(Do you want to use tools, or depend on them?)}

2. Môn học nào đang rèn cho em cách nghĩ chứ không chỉ cách trả lời? \textit{(Which subject is training your thinking, not just your answers?)}
\end{tuhocnhanh}
\ngancach

\section{Em sợ chọn sai nghề rồi mắc kẹt cả đời}
\begin{truyen}{Em sợ chọn sai nghề rồi mắc kẹt cả đời}{Classroom Talk}
\chuhoa{M}{ột bạn thở dài: ``Nghe người lớn hỏi chọn ngành là em mệt. Em sợ chọn sai một phát rồi kẹt luôn cả đời.''}

Thầy nói: ``Sợ vậy là bình thường. Nhưng đời không chỉ có một nút bấm. Chọn nghề quan trọng, nhưng khả năng học lại và chỉnh hướng còn quan trọng hơn.''

Bạn ấy hỏi: ``Vậy sai vẫn cứu được hả thầy?''

Thầy gật đầu: ``Sai có giá của nó. Nhưng nhiều người không chết vì chọn sai. Họ mòn vì chọn mà không chịu học tiếp.''

\textit{(A wrong choice matters, but adaptability often matters more than the illusion of a flawless first decision.)}
\end{truyen}

\begin{baihoc}
Không ai đảm bảo chọn đúng tuyệt đối. Điều làm mình đỡ kẹt là năng lực đổi hướng khi cần.
\textit{(No one can guarantee a perfect first choice. The real protection is the ability to redirect.)}
\end{baihoc}

\begin{ghinhoanh}
\textbf{Short English to remember:}

No first choice is perfect.

Adaptability is part of career strength.
\end{ghinhoanh}

\begin{tuhocnhanh}
1. Em đang sợ chọn sai hay sợ chịu trách nhiệm sau khi chọn? \textit{(Do you fear the wrong choice, or the responsibility after choosing?)}

2. Kỹ năng nào giúp em đổi hướng nếu cần? \textit{(What skills would help you change direction later?)}
\end{tuhocnhanh}
\ngancach

\section{Không có quen biết thì cố có ăn thua không?}
\begin{truyen}{Không có quen biết thì cố có ăn thua không?}{Classroom Talk}
\chuhoa{C}{ó bạn hỏi câu rất chát: ``Người ta có quan hệ, có chỗ đỡ đầu. Còn em không có ai. Vậy cố cho lắm có ăn thua không?''}

Thầy đáp: ``Quan hệ là thật. Nhưng không có năng lực mà chỉ có quan hệ thì đi không xa. Còn có năng lực mà không biết mở miệng, không biết làm việc, cũng tự chặn đường mình.''

Bạn ấy hỏi tiếp: ``Vậy em phải làm gì?''

Thầy nói: ``Vừa học cho chắc, vừa tập cách cư xử, làm việc, giao tiếp cho đàng hoàng. Đó cũng là một dạng vốn.''

\textit{(Connections matter, but so do competence and the social skills that help others trust your work.)}
\end{truyen}

\begin{baihoc}
Đừng ngây thơ nghĩ chỉ cần giỏi là đủ. Cũng đừng cay cú tới mức quên xây giá trị của mình.
\textit{(Do not naively think skill alone solves everything. But do not become so bitter that you stop building your own value.)}
\end{baihoc}

\begin{ghinhoanh}
\textbf{Short English to remember:}

Connections matter.

So do trust and competence.
\end{ghinhoanh}

\begin{tuhocnhanh}
1. Em đang thiếu năng lực, thiếu cơ hội, hay thiếu cách kết nối? \textit{(Are you lacking skill, opportunity, or the ability to connect?)}

2. Em có đang bỏ quên kỹ năng làm việc với người khác không? \textit{(Are you neglecting the skill of working with others?)}
\end{tuhocnhanh}
\ngancach